\documentclass[11pt]{beamer}
\usepackage[utf8]{inputenc}
\usepackage[T1]{fontenc}
\usepackage{lmodern}
\usetheme{Hannover}
\usepackage{amsmath}
\usepackage{amsthm}
\usepackage{graphicx} 
\begin{document}
	\author{Shiyu Dong}
	\title{The symmetric Group}
	%\subtitle{}
	%\logo{}
	%\institute{}
	%\date{}
	%\subject{}
	%\setbeamercovered{transparent}
	%\setbeamertemplate{navigation symbols}{}

	\begin{frame}[plain]
		\maketitle
	\end{frame}
\begin{frame}{outline}
	\tableofcontents
\end{frame}
		\section{Permutation}
	\begin{frame}{Permutation}
		\begin{equation}
			P=\begin{pmatrix}
				1&2&\cdots&n \\
				p_1&p_2&\cdots & p_n
			\end{pmatrix}
		\end{equation}
			\begin{equation}
		Q=\begin{pmatrix}
			1&2&\cdots&n \\
			q_1&q_2&\cdots & q_n
		\end{pmatrix}
	\end{equation}	
	\begin{align}
		PQ&=\begin{pmatrix}
			1&2&\cdots&n \\
			p_1&p_2&\cdots & p_n
		\end{pmatrix}\begin{pmatrix}
			1&2&\cdots&n \\
			q_1&q_2&\cdots & q_n
		\end{pmatrix} \\
		&=\begin{pmatrix}
				q_1&q_2&\cdots & q_n \\
			r_1&r_2&\cdots & r_n
		\end{pmatrix}\begin{pmatrix}
			1&2&\cdots&n \\
			q_1&q_2&\cdots & q_n
		\end{pmatrix} \\
		&=\begin{pmatrix}
			1&2&\cdots&n \\
			r_1&r_2&\cdots & r_n
		\end{pmatrix}
	\end{align}
		
	\end{frame}
\begin{frame}{cycle}
	\textit{m-cycle}   
	\begin{equation}
		\begin{pmatrix}
			1&2&\cdots & m-1 &m \\
			2&3&\cdots &m &1 
		\end{pmatrix}\equiv 
			\begin{pmatrix}
			1&2&\cdots & m-1 &m 
		\end{pmatrix}
	\end{equation}
	2-cycle is called \textit{transpositon}
	\begin{equation}
		\begin{pmatrix}
			1 &2 
		\end{pmatrix} \quad \cdots \quad 	\begin{pmatrix}
		1 &i
		\end{pmatrix} \quad \begin{pmatrix}
		i &j
		\end{pmatrix}
	\end{equation}
\end{frame}
	\begin{frame}{Parity}
		\begin{equation}
		\begin{split}
		\Delta(x_1,x_2,\cdots,x_n)=(x_1-x_2)(x_1-x_3)\cdots(x_1-x_n)\\
		\times(x_2-x_3)\cdots(x_2-x_n) \cdots \\
		\times(x_{n-1}-x_n)
	\end{split}		
		\end{equation}
\begin{equation}
	P \Delta=(-1)^{\sigma_p}\Delta
\end{equation}
\end{frame}
\begin{frame}{Decomposition}

		Every permutation can be expressed as a pdoduct of cycles 
	Example
	\begin{equation}
		\begin{pmatrix}
			1 &2&3&4&5&6\\
			2&4&5&1&3&6
		\end{pmatrix}\equiv \begin{pmatrix}
		1 &2&4 
		\end{pmatrix} \begin{pmatrix}
	3 &5 
		\end{pmatrix}
		\begin{pmatrix}
			6
		\end{pmatrix}
	\end{equation}
	Every permutation can be expressed as a product of transposition \\
	A m-cycle can be decomposed into product of m-1 transposition
\end{frame}
\begin{frame}{Decomposition}
	\begin{equation}
		\begin{pmatrix}
			1&2&3&\cdots& m-1&m\\
				1&2&3&\cdots& m-1&m
		\end{pmatrix}
	\end{equation}
	\begin{equation}
(1m)				\begin{pmatrix}
			1&2&\cdots& m-1&m\\
			1&2&\cdots& m-1&m
		\end{pmatrix}=	\begin{pmatrix}
		1&2&\cdots& m-1&m\\
		m&2&\cdots& m-1&1
		\end{pmatrix}
	\end{equation}
	\begin{equation}
(1m-1)	(1m)				\begin{pmatrix}
		1&2&\cdots& m-1&m\\
		1&2&\cdots& m-1&m
	\end{pmatrix}=	\begin{pmatrix}
		1&2&\cdots& m-1&m\\
		m-1&2&\cdots& m&1
	\end{pmatrix}
\end{equation}
\end{frame}
		\section{Partitions And conjugate classes}
\begin{frame}{Partition and conjugate class}
		\begin{equation}
		Q=TPT^{-1}
	\end{equation}
\begin{theorem}
	Two permutation in $S_n$ are conjugate if and only if they have the same cycle structure(they have the same number of cycles of equal length)
\end{theorem}
\end{frame}
\begin{frame}{Proof}
	Conjugate $\implies $same cycle structure
	\begin{equation}
		Q=TPT^{-1}=\begin{pmatrix}
			1&2\cdots&m\\
			t_1&t_2\cdots&t_m
		\end{pmatrix}\begin{pmatrix}
			1&2\cdots&m\\
			p_1&p_2\cdots&p_m
		\end{pmatrix}\begin{pmatrix}
		t_1&t_2\cdots&t_m\\
		1&2\cdots&m
		\end{pmatrix}
	\end{equation}
	\begin{equation}
	=\begin{pmatrix}
		t_1&t_2&\cdots & t_m\\
		p_{t_1}&p_{t_2}\cdots &p_{t_m}
	\end{pmatrix}
	\end{equation}
Relabeling the index.
\end{frame}
\begin{frame}{Cycle structure }
	\begin{equation}
		(\nu)=(1^{\nu_1}2^{\nu_2}\cdots m^{\nu_m})
	\end{equation}
	\begin{equation}
		\nu_1+2\nu_2+\cdots +m\nu_m=m
	\end{equation}
	\begin{equation}
		\lambda_1=\nu_1+\nu_2+\cdots +\nu_m
	\end{equation}
		\begin{equation}
		\lambda_k=\nu_k+\nu_2+\cdots +\nu_m
	\end{equation}
	\begin{equation}
		\lambda_1+\lambda_2+\cdots \lambda_m=m
	\end{equation}
	\begin{equation}
		\nu_k=\lambda_k-\lambda_{k+1}
	\end{equation}
	partition of integer $n$:
	\begin{equation*}
		\lambda_1+\lambda_2+\cdots+\lambda_n=n \quad \lambda_1\ge \lambda_2\cdots \lambda_n
	\end{equation*}
\end{frame}
		\section{Young Tableaux}
\begin{frame}{Symmetry and Antisymmety}
	% TODO: 
	\begin{figure}
		\centering
		\includegraphics[width=0.7\linewidth]{rep}
		\caption{}
		\label{fig:rep}
	\end{figure}
	
\end{frame}
		\section{Antisymmetrizer}
	\begin{frame}
			\begin{equation}
		 	\mathcal{\hat A}^{(n)}(1,2,\cdots,n)=\frac{1}{n!}\sum_P(-1)^{\sigma_p}\hat P
		\end{equation}
			\begin{equation}
			\hat Q\mathcal{\hat A}^{(n)} =(-1)^{\sigma_Q}\mathcal{\hat A}^{(n)}
		\end{equation}
		\begin{equation}
			\mathcal{\hat A}^{(n)}\mathcal{\hat A}^{(n)}=\mathcal{\hat A}^{(n)}
		\end{equation}
Proof
\begin{align}
	\hat Q\mathcal{\hat A}^{(n)}&=\frac{1}{n!}\sum_P (-1)^{\sigma_p}QP \\
	&=\frac{1}{n!}\sum_p (-1)^{\sigma_{Q}} (-1)^{\sigma_{QP}}QP  \\
	&= (-1)^{\sigma_{Q}}\frac{1}{n!}\sum_{qp} (-1)^{\sigma_{QP}}QP \\
	&= (-1)^{\sigma_{Q}}\mathcal{\hat A}^{(n)}
\end{align}
	\end{frame}
	\begin{frame}
	\begin{align}
		\mathcal{\hat A}^{(n)}\mathcal{\hat A}^{(n)}&=(\frac{1}{n!}\sum_{Q}(-1)^{\sigma_Q}Q)(\frac{1}{n!}\sum_{P}(-1)^{\sigma_P}P) \\
		&=\frac{1}{n!}\sum_{Q}(-1)^{\sigma_Q} (-1)^{\sigma_Q} 	\mathcal{\hat A}^{(n)}\\
		&=\hat A^{(n)}
	\end{align}
	\end{frame}
\end{document}